\documentclass[12pt]{article}
\usepackage{amsmath,amssymb,amsthm}
\usepackage{mathtools}
\usepackage{hyperref}
\usepackage{geometry}
\usepackage{booktabs}
\geometry{margin=1in}
\hypersetup{colorlinks=true,linkcolor=blue,citecolor=blue}

\newtheorem{theorem}{Theorem}[section]
\newtheorem{lemma}[theorem]{Lemma}
\newtheorem{corollary}[theorem]{Corollary}
\newtheorem{proposition}[theorem]{Proposition}
\theoremstyle{definition}
\newtheorem{definition}[theorem]{Definition}
\newtheorem{remark}[theorem]{Remark}
\newtheorem{openproblem}[theorem]{Open Problem}
\newtheorem{observation}[theorem]{Computational Observation}

\newcommand{\ZZ}{\mathbb{Z}}
\newcommand{\OK}{\Omega_K}
\newcommand{\vt}{v_2}
\newcommand{\Pk}{P_K}
\newcommand{\Pkm}{P_{K,m}}
\newcommand{\PKone}{P_{K,1}}
\newcommand{\pik}{\pi_K}
\newcommand{\TV}{\mathrm{TV}}

\title{Spectral Amplification by One-Bit Inflow\\
in Collatz Dynamics\\[0.5em]
{\large Paper C: Numerical Phase Diagram and an Open Gap Problem}}

\author{Hiroki Kamanoi}
\date{Draft --- May 2026}

\begin{document}
\maketitle

\begin{abstract}
We introduce and study a family of Markov kernels $\Pkm$ on odd residues modulo $2^K$, parametrized by the number of high bits $m$ averaged over. At $m=0$, $\Pkm$ reduces to the deterministic Collatz map $F_K$ with spectral gap zero. At $m=1$, the one-bit kernel $\PKone$ jumps to spectral gap $\delta_{K,1} \approx 0.29$, independently of $K$. This \emph{one-bit spectral amplification}
\[
  \delta_{K,0} = 0 \quad\xrightarrow{\;\text{1 bit}\;}\quad \delta_{K,1} \approx 0.29
\]
is the central phenomenon of this paper.

We record four computational observations: (i) spectral positivity $\delta_{K,1} \geq 0.27$ for $K = 4, \ldots, 13$; (ii) monotone increase of $\delta_{K,m}$ in $m$; (iii) the one-bit jump is the largest single increment; and (iv) saturation near $\delta_{K,K} \approx 0.74$.

We formulate the Open Gap Problem: $\inf_{K \geq 2} \delta_{K,1} > 0$. We prove that three standard techniques---Cheeger inequalities, Weyl perturbation, and Sinkhorn balancing---all fail to resolve this problem, and identify the core obstruction: the \emph{non-normal arithmetic geometry} of $\PKone$, whose spectral decomposition is governed by the global Collatz orbit structure rather than any local valuation invariant.
\end{abstract}

\tableofcontents\bigskip

%% ============================================================
\section{Introduction}
%% ============================================================

This is the third paper in a series studying Collatz dynamics through the lens of finite Markov chain theory. Paper B \cite{KamanoiB} established the exact TV mixing geometry of the true 2-adic Haar kernel $P_K$: the support of $P_K^t(a,\cdot)$ forms an exact arithmetic progression, and $\mathrm{TV}(P_K^t, \pi_K) = \max(0, 1-2^{S_t-(K-1)})$, yielding mixing time $T_{\mathrm{mix}}(K) = K-1$.

Paper B also revealed a fundamental separation:
\[
\underbrace{T_{\mathrm{mix}}(K) = K-1}_{\text{support mixing (TV)}} \quad\text{vs}\quad \underbrace{\delta_K \approx 2^{-(K-1)}}_{\text{spectral mixing (correlation)}}.
\]
The true 2-adic kernel has exponentially small spectral gap, meaning that while the support becomes uniform in $K-1$ steps, correlations between arithmetic functions of the orbit decay exponentially slowly.

The present paper asks: \emph{what happens when a finite amount of high-bit uncertainty is injected?} We show that even a single bit of high-bit uncertainty ($m=1$) jumps the spectral gap from $\approx 2^{-(K-1)}$ to $\approx 0.29$---a dramatic amplification that is \emph{uniform in $K$}.

\subsection*{A counterintuitive phenomenon}

The following comparison between Paper B and Paper C deserves emphasis, as it runs counter to naive intuition:
\[
  \underbrace{\delta_K^{\mathrm{Haar}} \approx 2^{-(K-1)} \to 0}_{\text{full 2-adic randomization}} \qquad\text{vs}\qquad \underbrace{\delta_{K,1} \approx 0.29}_{\text{one bit of structured inflow}}.
\]
One might expect that the fully Haar-random kernel---averaging over all 2-adic lifts---would mix faster than a single-bit perturbation. In fact, the opposite holds.

The mechanism is the following. The Haar kernel $P_K^{\mathrm{Haar}}$ spreads probability mass over $\Omega_K$ in a globally symmetric way, but in doing so it preserves long-range arithmetic Fourier coherence: certain arithmetic modes decay only at rate $2^{-(K-1)}$ per step because the Haar average cannot break them. The lone refresh point $a^*$ (the unique state with $\vt(3a^*+1) = K$) is responsible for all spectral mixing in $P_K^{\mathrm{Haar}}$, and it fires with probability only $1/n_K$ under the stationary measure.

In contrast, the one-bit kernel $\PKone$ acts asymmetrically: each state $a$ is displaced not to a globally uniform target, but to one of two specific arithmetic points $F_K(a)$ and $F_K(a) + 3\cdot 2^{K-v_a}$. The displacement $3\cdot 2^{K-v_a}$ is large (approximately $3n_K/2$ for Type A states), and this structured long-range jump immediately destroys low-bit arithmetic coherence at every state, not just at $a^*$. The result is a macroscopic spectral gap that is uniform in $K$.

In short: \emph{structured finite inflow beats full Haar randomization} in the spectral sense, because it disrupts arithmetic coherence globally rather than waiting for a single refresh event.

\subsection*{Main contributions}

\begin{enumerate}
\item Definition of the high-bit averaged kernel $\Pkm$ (Section~\ref{sec:def}).
\item Four computational observations establishing the spectral phase diagram, with the one-bit jump $\delta_{K,0} = 0 \to \delta_{K,1} \approx 0.29$ as the central finding (Section~\ref{sec:numerics}).
\item The Open Gap Problem: $\inf_K \delta_{K,1} > 0$ (Section~\ref{sec:open}).
\item A systematic negative result: three standard proof strategies fail, with identified failure mechanisms (Section~\ref{sec:negative}).
\item Heuristic connection to orbitwise decorrelation and the Collatz conjecture (Section~\ref{sec:collatz}).
\end{enumerate}

%% ============================================================
\section{Setup and Definitions}
\label{sec:def}
%% ============================================================

\begin{definition}[High-bit averaged kernel]
\label{def:Pkm}
For $K \geq 2$ and $m \geq 0$, the \emph{$m$-bit high-bit kernel} is
\[
  \Pkm(a,b) = \frac{1}{2^m} \#\!\left\{u \in \{0,\ldots,2^m-1\} : T(a + u \cdot 2^K) \equiv b \pmod{2^K}\right\},
\]
where $T(n) = (3n+1)/2^{\vt(3n+1)}$ is the Syracuse map and $a, b \in \OK = \{1, 3, \ldots, 2^K-1\}$.
\end{definition}

The kernel $\Pkm$ averages the Collatz image over $2^m$ lifts of $a$ to $a + u \cdot 2^K$. At $m=0$, this is the deterministic map $F_K(a) = T(a) \bmod 2^K$. As $m \to \infty$ (with the arithmetic average), the kernel approaches but does not converge to the true 2-adic Haar kernel $P_K^{\mathrm{Haar}}$.

\begin{definition}[Spectral gap]
The \emph{spectral gap} of $\Pkm$ is
\[
  \delta_{K,m} = 1 - |\lambda_2(\Pkm)|,
\]
where $\lambda_1 = 1 \geq |\lambda_2| \geq \ldots$ are the eigenvalues of $\Pkm$ ordered by absolute value.
\end{definition}

\begin{lemma}[V-stability]
\label{lem:vstab}
For odd $a \in \OK$ with $\vt(3a+1) = v_a < K$, and any $u \geq 0$:
\[
  T(a + u \cdot 2^K) \equiv T(a) + 3u \cdot 2^{K - v_a} \pmod{2^K}.
\]
In particular, the valuation $\vt(3(a+u\cdot 2^K)+1) = v_a$ is independent of $u$.
\end{lemma}

\begin{proof}
$3(a+u\cdot 2^K)+1 = (3a+1)+3u\cdot 2^K$. Since $v_a = \vt(3a+1) < K$ and $\vt(3u\cdot 2^K) \geq K > v_a$, the 2-adic valuation of the sum equals $v_a$. The stated congruence follows by dividing by $2^{v_a}$.
\end{proof}

\begin{corollary}
For $v_a < K$, the one-bit kernel satisfies
\[
  \PKone(a, \cdot) = \tfrac{1}{2}\,\delta_{F_K(a)} + \tfrac{1}{2}\,\delta_{F_K(a)+3\cdot 2^{K-v_a}}.
\]
\end{corollary}

\subsection{Basic properties of $\PKone$}

\begin{proposition}[Non-ideal properties]
\label{prop:nonideal}
The kernel $\PKone$ satisfies:
\begin{enumerate}
\item \textbf{Row-stochastic}: $\sum_b \PKone(a,b) = 1$ for all $a$.
\item \textbf{Not doubly stochastic}: column sums range in $[0.5,\, O(K)]$.
\item \textbf{Not reversible with respect to its stationary measure}: detailed balance $\pi_{K,1}(a)\PKone(a,b) = \pi_{K,1}(b)\PKone(b,a)$ fails for generic pairs $(a,b)$.
\item \textbf{Non-uniform stationary}: the stationary distribution $\pi_{K,1}$ is not uniform; it concentrates mass at $a=1$.
\item \textbf{Non-local eigenvectors}: the second eigenvector has no apparent structure in terms of $\vt(3a+1)$.
\end{enumerate}
\end{proposition}

These properties collectively make $\PKone$ analytically harder to control than standard reversible chains, and motivate the formulation of an open problem (Section~\ref{sec:open}).

%% ============================================================
\section{Numerical Phase Diagram}
\label{sec:numerics}
%% ============================================================

We computed $\delta_{K,m}$ for $K = 4, \ldots, 13$ and $m = 0, \ldots, K$ using exact eigenvalue computation via the QR algorithm.

\subsection{Spectral positivity at $m=1$}

\begin{observation}[Spectral Positivity]
\label{obs:positivity}
For $K = 4, 5, \ldots, 13$, the spectral gap satisfies
\[
  \delta_{K,1} \in [0.272,\; 0.355].
\]
In particular, $\inf_{K=4}^{13} \delta_{K,1} \geq 0.272 > 0$.
\end{observation}

\begin{center}
\begin{tabular}{@{}crrrr@{}}
\toprule
$K$ & $n_K$ & $\lambda_1$ & $|\lambda_2|$ & $\delta_{K,1}$ \\
\midrule
 4 &    8 & 1.000 & 0.711 & 0.289 \\
 5 &   16 & 1.000 & 0.645 & 0.355 \\
 6 &   32 & 1.000 & 0.687 & 0.313 \\
 7 &   64 & 1.000 & 0.698 & 0.302 \\
 8 &  128 & 1.000 & 0.703 & 0.297 \\
 9 &  256 & 1.000 & 0.722 & 0.278 \\
10 &  512 & 1.000 & 0.724 & 0.276 \\
11 & 1024 & 1.000 & 0.710 & 0.290 \\
12 & 2048 & 1.000 & 0.728 & 0.272 \\
13 & 4096 & 1.000 & 0.719 & 0.281 \\
\bottomrule
\end{tabular}
\end{center}

\subsection{Phase monotonicity}

\begin{observation}[Phase Monotonicity]
\label{obs:mono}
For each $K \in \{6, 8, 10\}$, the map $m \mapsto \delta_{K,m}$ is monotone non-decreasing on $\{0, 1, \ldots, K\}$.
\end{observation}

\begin{observation}[One-Bit Jump]
\label{obs:jump}
The largest marginal gain per bit occurs at $m = 1$:
\[
  \delta_{K,1} - \delta_{K,0} = \delta_{K,1} - 0 = \delta_{K,1} \approx 0.29.
\]
Subsequent increments $\delta_{K,m+1} - \delta_{K,m}$ are smaller and decreasing.
\end{observation}

\begin{observation}[Saturation]
\label{obs:sat}
The spectral gap saturates near $m = K$:
$\delta_{K,K} \approx 0.74$ (independent of $K$), and the final increment $\delta_{K,K} - \delta_{K,K-1} \leq 0.022$.
\end{observation}

\subsection{Phase diagram summary}

\begin{center}
\begin{tabular}{@{}ccccccccccc@{}}
\toprule
$K \backslash m$ & 0 & 1 & 2 & 3 & 4 & 5 & 6 & 7 & 8 & $K$ \\
\midrule
6  & 0 & .313 & .428 & .571 & .610 & .701 & .723 & -- & -- & .723 \\
8  & 0 & .297 & .421 & .552 & .598 & .661 & .686 & .724 & .745 & .745 \\
10 & 0 & .276 & .403 & .482 & .544 & .578 & .623 & .661 & .681 & .730 \\
\bottomrule
\end{tabular}
\end{center}

The regime structure is:
\begin{itemize}
\item \textbf{Regime I} ($m=0$): $\delta_{K,0} = 0$ (deterministic).
\item \textbf{Regime II} ($1 \leq m \leq K$): $\delta_{K,m} \approx$ constant in $K$ for fixed $m/K$.
\item \textbf{Regime III} (true Haar, $m \to \infty$ in 2-adic sense): $\delta_K^{\mathrm{Haar}} \approx 2^{-(K-1)} \to 0$.
\end{itemize}

%% ============================================================
\section{The Open Gap Problem}
\label{sec:open}
%% ============================================================

The numerical evidence motivates the following central open problem.

\begin{openproblem}[Open Gap Problem]
\label{prob:gap}
Prove or disprove:
\[
  \boxed{\inf_{K \geq 2} \delta_{K,1} > 0.}
\]
\end{openproblem}

\begin{remark}
The physical interpretation is: \emph{a single bit of high-bit uncertainty injects a uniformly positive amount of decorrelation into the Collatz dynamics at every scale $K$}. If Open Problem~\ref{prob:gap} is true, it means that the deterministic Collatz map, once perturbed by one bit of arithmetic ambiguity, cannot sustain long-range spectral coherence---a necessary condition for orbitwise decorrelation.
\end{remark}

\begin{remark}[Equivalent formulations]
Open Problem~\ref{prob:gap} is equivalent to each of:
\begin{enumerate}
\item $\liminf_{K\to\infty} \delta_{K,1} > 0$.
\item $\sup_{K} |\lambda_2(\PKone)| < 1$.
\item There exists $t_0 < \infty$ such that $\sup_K \|\PKone^{t_0} - \pi_{K,1}\|_{\mathrm{op}} < 1$.
\end{enumerate}
\end{remark}

%% ============================================================
\section{Why Standard Approaches Fail}
\label{sec:negative}
%% ============================================================

We prove that three standard techniques are insufficient for resolving Open Problem~\ref{prob:gap}, and identify the core obstruction.

\subsection{Cheeger inequality fails}

The Cheeger inequality for reversible chains states $\delta \geq h^2/2$ where $h$ is the conductance. This requires the chain to be reversible (detailed balance) with uniform stationary distribution.

\begin{proposition}
$\PKone$ is not reversible with respect to its stationary measure $\pi_{K,1}$. Moreover, numerically $\delta(\tilde{P}) \to 0$ as $K \to \infty$ for the additive symmetrization $\tilde{P} = \frac{1}{2}(\PKone + \PKone^T)$.
\end{proposition}

One may attempt to use the additive symmetrization $\tilde{P} = \frac{1}{2}(\PKone + \PKone^T)$. However:

\begin{proposition}
The spectral gap of the symmetrization satisfies $\delta(\tilde{P}) \to 0$ as $K \to \infty$.
\end{proposition}

This means Cheeger applied to $\tilde{P}$ provides no useful lower bound on $\delta_{K,1}$.

\subsection{Weyl perturbation fails}

Write $\PKone = \frac{1}{2}(P^{\mathrm{det}} + P^{\mathrm{flip}})$, where $P^{\mathrm{det}}(a,\cdot) = \delta_{F_K(a)}$ and $P^{\mathrm{flip}}(a,\cdot) = \delta_{T(a+2^K) \bmod 2^K}$.

Weyl's inequality gives $\lambda_j(A+B) \leq \lambda_j(A) + \|B\|_{\mathrm{op}}$. Here $\|P^{\mathrm{flip}}\|_{\mathrm{op}} = 1$ and $\lambda_2(P^{\mathrm{det}}) = 1$ (or 0 for nilpotent $F_K$), so the bound yields $\delta_{K,1} \geq -1$, which is trivial.

The fundamental reason is that both summands have operator norm $1$, so their mixture can have eigenvalue anywhere in $[0,1]$. Weyl does not resolve the problem.

\subsection{Sinkhorn balancing fails}

One may hope to balance $\PKone$ to a doubly stochastic matrix $\tilde{P}$ via Sinkhorn-Knopp iteration, then apply Cheeger. However:

\begin{proposition}
The support graph of $\PKone$ has in-degree distribution ranging from $1$ to $O(K)$. Nodes with in-degree $1$ prevent convergence of Sinkhorn balancing to a uniform doubly stochastic matrix. Numerically, $\delta(\tilde{P}_{K,1}) \to 0$ as $K \to \infty$.
\end{proposition}

\subsection{The core obstruction}

\begin{remark}[Core obstruction]
All three approaches fail for the same fundamental reason: the spectral decomposition of $\PKone$ is governed by the \emph{global Collatz orbit structure} rather than any local valuation invariant. Specifically:
\begin{itemize}
\item The second eigenvector of $\PKone$ has no simple structure in terms of $\vt(3a+1)$.
\item $\PKone$ is \emph{non-normal} ($\PKone \PKone^T \neq \PKone^T \PKone$), so its eigenvalues cannot be read off from any local quantity.
\item The stationary measure $\pi_{K,1}$ concentrates near $a=1$ (the Collatz fixed point), reflecting the global convergence of $F_K$ rather than any arithmetic symmetry.
\end{itemize}

The obstruction is not a lack of expansion, but the \emph{non-normal arithmetic geometry} of the Collatz transition operator: standard mixing theory requires either geometric expansion (Cheeger) or norm control (Weyl), and $\PKone$ satisfies neither condition in a $K$-uniform way.
\end{remark}

%% ============================================================
\section{Heuristic Connection to Orbitwise Decorrelation}
\label{sec:collatz}
%% ============================================================

\emph{The following discussion is heuristic and intended to motivate the relevance of Open Problem~\ref{prob:gap} to orbitwise Collatz dynamics. No claims here are asserted as theorems.}

We explain why Open Problem~\ref{prob:gap} is connected to the Collatz conjecture.

\subsection{The orbitwise interpretation}

Let $n_0$ be a positive integer and $n_t = T^t(n_0)$ its Collatz orbit. At time $t$, the residue $n_t \bmod 2^K$ evolves under the deterministic kernel $P^{\mathrm{det}}$. However, the high bits $\lfloor n_t / 2^K \rfloor$ evolve non-trivially and feed information back into the low bits.

By Theorem ETV (Paper B), after $K-1$ steps, the conditional distribution of $n_t \bmod 2^K$ given the first $K$ bits of $n_0$ becomes exactly uniform. This is the \emph{TV mixing} phenomenon.

However, TV mixing does not imply spectral decorrelation. For spectral decorrelation, one needs:
\[
  \left|\mathbb{E}[f(n_t \bmod 2^K) \mid n_0 \bmod 2^K = a] - \int f\,d\pi_K\right| \leq C \cdot (1-\delta_K)^t \cdot \|f\|_\infty.
\]

\subsection{Effective high-bit inflow}

In the actual Collatz orbit, the high bits $\lfloor n_t/2^K \rfloor$ are deterministic but grow rapidly. After $t$ steps, effectively $m \approx t$ bits of new high-bit information have flowed into the low $K$ bits.

Thus the relevant spectral gap for the orbit at time $t$ is approximately $\delta_{K,m}$ with $m \approx t$. By Computational Observation~\ref{obs:mono}, $\delta_{K,m}$ is increasing in $m$. By Computational Observation~\ref{obs:jump}, the most critical step is $m=1$ (the first bit of inflow).

\subsection{Consequence of the Open Gap Problem}

If Open Problem~\ref{prob:gap} holds with constant $c > 0$, then for any $K$ and any $t \geq 1$:
\[
  \left|\mathbb{E}[f(n_t \bmod 2^K) \mid n_0] - \int f\,d\pi_{K,1}\right| \leq (1-c)^{t-1} \cdot \|f\|_\infty.
\]
This would establish that \emph{after one step of high-bit inflow, arithmetic correlations decay at a uniform exponential rate, independent of the scale $K$}.

A divergent Collatz orbit would need to maintain arithmetic coherence across all scales $K$ simultaneously, but if each scale decorrelates at rate $c$ per step of high-bit inflow, the total coherence would decay geometrically. This is a strong obstruction to divergence.

%% ============================================================
\section{Summary}
%% ============================================================

\begin{center}
\begin{tabular}{@{}lll@{}}
\toprule
Result & Status & Section \\
\midrule
Definition of $\Pkm$ & Exact & \S\ref{sec:def} \\
Lemma V-stability & Proved & \S\ref{sec:def} \\
Basic properties of $\PKone$ & Proved & \S\ref{sec:def} \\
Obs.\ 1: Spectral Positivity $\delta_{K,1} \geq 0.27$ & Computational & \S\ref{sec:numerics} \\
Obs.\ 2: Phase Monotonicity in $m$ & Computational & \S\ref{sec:numerics} \\
Obs.\ 3: One-Bit Jump is largest increment & Computational & \S\ref{sec:numerics} \\
Obs.\ 4: Saturation at $\delta_{K,K} \approx 0.74$ & Computational & \S\ref{sec:numerics} \\
Open Gap Problem: $\inf_K \delta_{K,1} > 0$ & \textbf{Open} & \S\ref{sec:open} \\
Cheeger fails (non-reversible, non-uniform $\pi$) & Proved & \S\ref{sec:negative} \\
Weyl fails ($\|P^{\mathrm{flip}}\|_{\mathrm{op}} = 1$) & Proved & \S\ref{sec:negative} \\
Sinkhorn fails (in-degree imbalance) & Proved & \S\ref{sec:negative} \\
Core obstruction: non-normal arithmetic geometry & Identified & \S\ref{sec:negative} \\
Orbitwise decorrelation connection & Heuristic & \S\ref{sec:collatz} \\
\bottomrule
\end{tabular}
\end{center}

\begin{thebibliography}{99}

\bibitem{KamanoiA}
H. Kamanoi,
``Exact Operator Structure and Finite Information Memory in Accelerated Collatz Dynamics''
(Paper A), 2025.

\bibitem{KamanoiB}
H. Kamanoi,
``Exact Finite-Level Mixing Geometry in Collatz Dynamics''
(Paper B), 2026.

\bibitem{Lagarias2010}
J.~C. Lagarias (ed.),
\textit{The Ultimate Challenge: The $3x+1$ Problem},
AMS, 2010.

\bibitem{Levin2017}
D.~A. Levin, Y. Peres, E.~L. Wilmer,
\textit{Markov Chains and Mixing Times}, 2nd ed., AMS, 2017.

\bibitem{Sinkhorn1967}
R. Sinkhorn and P. Knopp,
``Concerning nonnegative matrices and doubly stochastic matrices,''
\textit{Pacific J. Math.} \textbf{21} (1967), 343--348.

\end{thebibliography}

\end{document}
